\section{D. Penerapan Nilai-Nilai Pancasila}

Pancasila adalah pandangan hidup bangsa Indonesia. Sebelum terbentuknya Negara Indonesia yang diproklamasikan tanggal 17 Agustus 1945, Pancasila telah terlihat dan diterapkan dalam kehidupan sehari-hari yang digali dari nilai adat istiadat, kebudayaan, serta nilai religius.

\subsection{Wawasan Kewarganegaraan: Kitab Sutasoma}
\begin{itemize}
    \item Ditulis dalam bahasa Jawa Kuno oleh Mpu Tantular pada akhir abad ke-14 masa Kerajaan Majapahit.
    \item Menggambarkan toleransi beragama antara ajaran Buddha dan Siwa.
    \item Kutipan frase \textit{Bhinneka Tunggal Ika} terdapat pada pupuh 139 bait 5: \textit{"Rwaneka dhatu winuwus Buddha Wiswa, Bhinneki rakwa ring apan kena parwanosen, Mangka ng Jinatwa kalawan Siwatatwa tunggal, Bhinneka tunggal ika tan hana dharma mangrwa"}.
\end{itemize}

\subsection{Asal Mula Pancasila}
Menurut Kaelan (2020) mengutip Notonagoro, asal mula Pancasila dibedakan menjadi:

\subsubsection{Asal Mula Langsung}
\begin{enumerate}
    \item \textbf{Kausa Materialis (Asal Mula Bahan)}: Nilai-nilai Pancasila digali dari bangsa Indonesia sendiri berupa nilai adat istiadat, kebudayaan, serta nilai religius.
    \item \textbf{Kausa Formalis (Asal Mula Bentuk)}: Dirumuskan sejak usulan Ir. Sukarno pada sidang pertama BPUPK, dibahas oleh BPUPK dan Panitia Sembilan, hingga disahkan oleh PPKI pada 18 Agustus 1945.
    \item \textbf{Kausa Efisien (Asal Mula Karya)}: Asal mula yang menjadikan Pancasila dari calon dasar negara menjadi dasar negara yang sah oleh PPKI.
    \item \textbf{Kausa Finalis (Asal Mula Tujuan)}: Dirumuskan dan dibahas dalam persidangan oleh pendiri negara untuk ditetapkan sebagai dasar negara Indonesia Merdeka.
\end{enumerate}

\subsubsection{Asal Mula Tidak Langsung}
Nilai-nilai Pancasila yang terdapat dalam adat istiadat, kebudayaan, serta nilai agama bangsa Indonesia sebelum Proklamasi Kemerdekaan.

\subsubsection{Tataran Nilai Ideologi Pancasila}
\begin{itemize}
    \item \textbf{Nilai Dasar}: Nilai fundamental dan universal (Ketuhanan, Kemanusiaan, Persatuan, Kerakyatan, Keadilan).
    \item \textbf{Nilai Instrumental}: Penjabaran nilai dasar dalam bentuk peraturan perundang-undangan dan kebijakan.
    \item \textbf{Nilai Praksis}: Realisasi nilai-nilai dalam kehidupan bermasyarakat, berbangsa, dan bernegara.
\end{itemize}

\subsection{Wawasan Kewarganegaraan: Investasi Mental Karakter}
\begin{itemize}
    \item Gagasan Ir. Sukarno mengenai \textit{Nation and Character Building} melalui tiga fase revolusi bangsa:
    \begin{enumerate}
        \item \textit{Physical Revolution} (1945--1949)
        \item \textit{Survival} (1950--1955)
        \item \textit{Investment} (\textit{investment of human skill}, \textit{material investment}, dan \textit{mental character investment}).
    \end{enumerate}
\end{itemize}

\subsection{1. Pembudayaan Pancasila}
Menurut Yudi Latif (2020), pembudayaan Pancasila ditempuh melalui tiga ranah peradaban:
\begin{itemize}
    \item \textbf{Ranah Keyakinan}: Penghayatan Pancasila dalam kerangka tertib sosial.
    \item \textbf{Ranah Pengetahuan}: Pengetahuan tentang Pancasila sebagai kerangka paradigmatik (\textit{civic intelligence}).
    \item \textbf{Ranah Tindakan}: Komitmen mengamalkan Pancasila dalam tindakan dan kebajikan.
\end{itemize}

\subsection{Wawasan Kewarganegaraan: Ki Hajar Dewantara, Sebuah Inspirasi}
Tiga Prinsip Utama Kepemimpinan Ki Hajar Dewantara:
\begin{enumerate}
    \item \textbf{Ing Ngarsa Sung Tulada}: Di depan memberi teladan.
    \item \textbf{Ing Madya Mangun Karsa}: Di tengah membangun semangat/prakarsa.
    \item \textbf{Tut Wuri Handayani}: Di belakang memberikan dorongan moral.
\end{enumerate}

\subsection{2. Strategi Perubahan dan Pengamalan}
Lima langkah prioritas agenda transformasi menurut Yudi Latif (2020):
\begin{enumerate}
    \item Melakukan revitalisasi dan reaktualisasi pemahaman terhadap Pancasila.
    \item Mengembangkan kerukunan (inklusi sosial) di tengah masyarakat.
    \item Mendorong terwujudnya keadilan sosial melalui sistem ekonomi berbasis Pancasila.
    \item Menguatkan internalisasi nilai-nilai Pancasila ke dalam produk perundang-undangan.
    \item Menumbuhkan, mempromosikan, dan mengapresiasi keteladanan agen-agen kenegaraan dan kemasyarakatan.
\end{enumerate}

\subsubsection{Jalur Utama Peradaban Pancasila}
\begin{table}[h!]
\centering
\begin{tabular}{|l|l|}
\hline
\textbf{Jalur Perubahan} & \textbf{Tujuan Peradaban} \\ \hline
Jalur Pemahaman & Indonesia Cerdas Kewargaan \\ \hline
Jalur Kerukunan & Indonesia Bersatu \\ \hline
Jalur Keadilan & Indonesia Berbagi Sejahtera \\ \hline
Jalur Pelembagaan & Indonesia Tertata Terlembaga \\ \hline
Jalur Keteladanan & Indonesia Terpuji \\ \hline
\end{tabular}
\end{table}

\textit{Catatan Kelembagaan}: Pemerintah membentuk Badan Pembinaan Ideologi Pancasila (BPIP) pada 28 Februari 2018 berdasarkan Peraturan Presiden Nomor 7 Tahun 2018.

\subsection{Contoh Penerapan Sikap dan Perilaku Terhadap Nilai-Nilai Pancasila}

\subsubsection{a. Di Lingkungan Keluarga}
\begin{enumerate}
    \item Memecahkan masalah dengan berdiskusi bersama anggota keluarga.
    \item Menghormati segenap anggota keluarga.
    \item Kasih sayang antarsaudara.
    \item Membantu menyelesaikan pekerjaan di rumah.
\end{enumerate}

\subsubsection{b. Di Lingkungan Sekolah}
\begin{enumerate}
    \item Disiplin mematuhi tata tertib sekolah.
    \item Toleransi terhadap keberagaman antarsiswas.
    \item Membantu teman yang sedang terkena musibah.
    \item Melaporkan terjadinya kasus perundungan kepada guru.
\end{enumerate}

\subsubsection{c. Di Lingkungan Masyarakat}
\begin{enumerate}
    \item Rukun dengan sesama tetangga.
    \item Mengutamakan musyawarah dalam pengambilan keputusan bersama.
    \item Menghormati sesama warga.
    \item Tolong-menolong membantu warga yang terkena musibah.
\end{enumerate}

\subsubsection{d. Di Lingkungan Bangsa dan Negara}
\begin{enumerate}
    \item Menyampaikan pendapat atau aspirasi secara bertanggung jawab.
    \item Toleransi terhadap keberagaman.
    \item Memberikan bantuan pada saudara sebangsa yang sedang tertimpa bencana alam.
    \item Mematuhi hukum yang berlaku.
\end{enumerate}

\subsection{Wawasan Kewarganegaraan: Ribuan Anak Yatim Luar Panti Terima Santunan}
Pemerintah Kabupaten Semarang menyalurkan bantuan sosial/santunan kepada anak-anak yatim luar panti sebagai wujud komitmen jaminan kesejahteraan sosial warga.